\chapter{Appendix 附錄}
\label{ch:appendix}

本附錄整理撰寫論文時常見的排版與語法注意事項，供大家撰寫與校稿時參考，請記得看latex的語法。(此段落主要由ChatGPT撰寫整理 XD)

\section*{文字與標點}
\begin{itemize}
  \item 英文引號請使用成對反引號與直引號：``text'' → “text”。
  \item 破折號（dash）有三種：`-`（連字號，如 \textit{pre-trained}）、`--`（範圍，如 10--20）、`---`（句中強調，如 the result---as expected---was correct）。
  \item 省略號請使用 \verb|\ldots| 而非三個句點：A, B, C\ldots, Z。
  \item 數值與單位間需留不換行空白：10~km，30$^\circ$C。
  \item 中英文混排時，中文與英文之間宜留一半形空白，例如：中文 English 中文。
\end{itemize}

\section*{數學與符號}
\begin{itemize}
  \item 數學變數自動為斜體，單位或文字常數須使用 \verb|\mathrm{}| 或 \verb|\text{}|，如 $v=10~\mathrm{m/s}$。
  \item 三角函數、極值運算等應以數學命令表示：$\sin\theta$, $\max(x,y)$。
  \item 向量或矩陣可用粗體：$\mathbf{x}=[x_1,x_2]^\top$。
  \item 多行方程式請使用 \texttt{align} 環境，勿以多個 \texttt{equation} 拆寫。
\end{itemize}


\section*{圖表與引用}
\begin{itemize}
  \item 圖片標題（caption）置於圖下方
  \item 表格標題置於表上方。
  \item 參照圖表時請使用 \verb|\ref{}| 或 \verb|\autoref{}|，勿手動輸入編號。\\
        例：As shown in Fig.~\ref{fig:excel2latex}.
  \item 表格建議使用 \texttt{booktabs} 套件：\verb|\toprule|、\verb|\midrule|、\verb|\bottomrule|。
\end{itemize}

\section*{文獻與引用}
\begin{itemize}
  \item 文獻引用使用 \verb|\cite{key}| 或 \verb|\parencite{key}|。
  \item 多篇同時引用可寫為 \verb|\cite{a,b,c}| → [1–3]。
  \item 特殊大小寫（如 GPU, LaTeX）須以大括號保留：\verb|{GPU}|。
\end{itemize}
